% Compiler avec XeLaTeX ou LuaLaTeX
\documentclass[11pt,a4paper]{article}

\usepackage{amsmath,amssymb}
\usepackage{fontspec}
\usepackage{unicode-math}
\usepackage{polyglossia}
\setmainlanguage{french}
\usepackage{xcolor}
\usepackage{fancyhdr}
\usepackage[margin=2cm,headheight=16pt,headsep=22pt]{geometry}
\usepackage{enumitem}

% Police principale et police mathématique : Libertinus
\IfFontExistsTF{Libertinus Serif}
  {\setmainfont{Libertinus Serif}}
  {\IfFontExistsTF{Latin Modern Roman}{\setmainfont{Latin Modern Roman}}{\setmainfont{lmroman10-regular.otf}}}

\IfFontExistsTF{Libertinus Math}
  {\setmathfont{Libertinus Math}}
  {\IfFontExistsTF{Latin Modern Math}{\setmathfont{Latin Modern Math}}{\setmathfont{latinmodern-math.otf}}}

\linespread{1.07}
\setlength{\parindent}{0pt}
\AtBeginDocument{\setlength{\parindent}{0pt}}
\setlength{\parskip}{0.35em}

% En-tête
\pagestyle{fancy}
\fancyhf{}
\lhead{Mathématiques}
\rhead{Première Générale -- Spécialité}
% Pagination : « Page X / Y » centrée en pied de page
\usepackage{lastpage}
\cfoot{Page \thepage\ / \pageref{LastPage}}
\renewcommand{\headrulewidth}{0.5pt}
% Trait de séparation éloigné du texte de l'en-tête (modifier 5pt pour ajuster)
\renewcommand{\headrule}{\vskip 5pt\hrule width\headwidth height\headrulewidth \vskip-\headrulewidth}

% Configuration des listes
\setlist[enumerate,1]{itemsep=0.45em, topsep=0.3em, label=\textbf{\arabic*.}}
\setlist[enumerate,2]{itemsep=0.25em, topsep=0.2em, label=\textbf{\alph*.}}

% Commande exercice
\newcommand{\exo}[2]{%
  \vspace{1.5em}%
  \noindent{\large\textbf{Exercice #1. -- #2}}%
  \par\vspace{0.35em}%
  \hrule height 0.5pt%
  \vspace{0.7em}%
}

\newcommand{\partie}[1]{%
  \vspace{0.9em}%
  \noindent\textbf{#1}%
  \par\vspace{0.3em}%
}

\begin{document}

\begin{center}
  {\LARGE \textbf{Dérivation et applications}}\\[0.5em]
  {\large Première Spécialité Mathématiques}
\end{center}

% ---------------------- EXERCICE 1 ----------------------
\exo{1}{Étude d'une fonction polynôme du troisième degré et recherche de tangentes}

Soit $f$ la fonction définie et dérivable sur $\mathbb{R}$ par :
\[
f(x) = 2x^3 - 3x^2 - 12x + 5
\]
On note $\mathcal{C}_f$ sa courbe représentative dans un repère orthonormé $(O; \vec{i}, \vec{j})$.

\partie{Partie A -- Étude des variations de la fonction $f$}

\begin{enumerate}
  \item Justifier que $f$ est dérivable sur $\mathbb{R}$ et déterminer sa fonction dérivée $f'$.
  \item Étudier le signe du trinôme $f'(x)$ selon les valeurs du réel $x$.
  \item Dresser le tableau de variations complet de la fonction $f$ sur l'intervalle $[-3\,;\,4]$, en précisant les valeurs remarquables et les extremums locaux.
  \item Préciser le nombre de solutions de l'équation $f(x) = 0$ sur l'intervalle $[-3\,;\,4]$. Justifier.
\end{enumerate}

\partie{Partie B -- Tangentes remarquables à la courbe $\mathcal{C}_f$}

\begin{enumerate}
  \item Déterminer l'équation réduite de la tangente $T_0$ à la courbe $\mathcal{C}_f$ au point d'abscisse $x = 0$.
  \item Déterminer les abscisses des éventuels points de $\mathcal{C}_f$ où la tangente est parallèle à la droite d'équation $y = 24x - 7$.
  \item Existe-t-il des points de $\mathcal{C}_f$ où la tangente est horizontale ? Si oui, préciser leurs coordonnées.
\end{enumerate}

\end{document}
